\scanpage{4}
\mspage{1}
Structure at infinity of the $M_{g,\nu}$

1. \uline{Standard curves}.

Let $k$ be an algebraically closed field. A ``standard curve'' over $k$ is a\deleted{ \unsure{curve}} scheme $X$ over $k$ satisfying the following conditions:

a) $X$ quasi-projective, \inserted{every irreducible component is} of dimension $1$

b) Every point of $X$ is either smooth, or a ``quadratic point''(ordinary'') -- i.e. isomorphic (étale locally) to \unsure{a|\deleted{a}} the curve $\operatorname{Spec}(k[X,Y]/XY)$ at the point $0$.

It is known that one can give a unique ``completion'' $\widehat X$ of $X$, such that \inserted{$\widehat X$ is a proper scheme, that} $X$ is identified with a dense open subset of $\widehat X$, and that $\widehat X$ is smooth at the points of $\widehat X\setminus X=I$. Then $\widehat X$ is a projective curve, $I$ is \deleted{\unsure{contained|curve}} a finite subset of $\widehat X(k)$ contained in the open set of smooth points of $\widehat X$. The set $A$ of singular points of $X$ is identified with the analogous set for $\widehat X$
\[
(1)\qquad
\left\{\begin{aligned}
I&=\widehat X\setminus X\subset\widehat X^{\mathrm{lisse}}(k)
&&\text{``points at infinity'' of }X\\
A&=X\setminus X^{\mathrm{lisse}}=\widehat X\setminus\widehat X^{\mathrm{lisse}}\subset X(k)
&&\text{(singular points of }X\text{)}
\end{aligned}\right.
\]
Giving $X$ is equivalent to giving
\mspage{3}
